\scanpage{33}
\mspage{4}
en l'ens. des \unsure{nb\textsuperscript{s}} premiers qui sont distincts
des caractéristiques résiduelles qui intervien-
nent, vu que ce n'est que dans ce
contexte modifié que l'on peut prouver
les théorèmes qu'il faut \deleted{de trivialité},
cf Artin-Mazur ...).

On aura en particulier une suite
exacte d'homotopie \deleted{\unsure{longue}} (choisissant
un \deleted{\unsure{pt base}} ``pt base'' dans une série
convenable \deleted{\unsure{p}} des germes d'espace $V^*_{z,Z}$) :

\noindent (12)
\diagram{\begin{tikzcd}[column sep=small]
\cdots & \pi_i(V^*_{\deleted{\unsure{Z}}z,Z}) \arrow[r] & \pi_i(V^*_{Y,X}) \arrow[r] & \pi_i(Y) \arrow[r] & \pi_{i-1}(V^*_{z,Z}) \arrow[r] & \cdots
\end{tikzcd}},

qui nous intéresse surtout en basses
dimensions

\noindent (13)
\diagram{\begin{tikzcd}[column sep=small,row sep=large]
\pi_2(Y) \arrow[r] & \pi_1(V^*_{z,Z}) \arrow[r] & \pi_1(V^*_{Y,X}) \arrow[r] & \pi_1(Y) \arrow[r] & \pi_0(V^*_{z,Z}) \arrow[d] \\
&&&& \pi_0(V^*_{Y,X}) \arrow[d] \\
&&&& \pi_0(Y)
\arrow[draw=none,from=1-3,to=1-4,"{\underbrace{\hspace{3.5cm}}_{\text{qui intervient dans (10)}}}"',shift right=1.5em]
\end{tikzcd}}

Si on suppose p. ex. que $z$ ne disconnecte
pas $Z$ localement i.e. : \deleted{\unsure{donc}} la ``fibre $V^*_{z,Z}$ est connexe
(ou encore, que $Y$ ne disconnecte pas $X$
loc\textsuperscript{t}.), alors $\pi_1(V^*_{Y,X})$ est une \uline{extension}
de $\pi_1(Y)$ par un quotient de $\pi_1(V^*_{z,Z})$, \deleted{\unsure{et}}
\mspage{32}
